\section{Cryptographic Case Studies}
\label{sec:case-studies}

The abstractions are exercised by four construction families.  The point of
these examples is not the novelty of their reductions, but the fact that exact
execution, cost erasure, and complexity evidence use the same interfaces.
\Cref{tab:case-studies} summarizes the mechanized boundary.

\begin{table}[H]
  \centering
  \caption{Construction-specific evidence connected to the common
    infrastructure.}
  \label{tab:case-studies}
  \small
  \begin{tabular}{@{}>{\raggedright\arraybackslash}p{.11\linewidth}
      @{\hspace{.02\linewidth}}
      >{\raggedright\arraybackslash}p{.39\linewidth}
      @{\hspace{.02\linewidth}}
      >{\raggedright\arraybackslash}p{.38\linewidth}@{}}
    \toprule
    Family & Exact program and bound evidence & Cost-erased result \\
    \midrule
    OTP & setup, key generation, encryption, and decryption programs;
      parameter and family efficiency certificates & correctness and perfect
      one-time security \\
    DLog & setup, challenge, and sampling programs; parameter/family
      efficiency certificates & search problem and all-PPT assumption \\
    DDH & setup, real/random challenge and sample programs over a shared
      decisional cyclic action & distinguishing problem and all-PPT assumption \\
    ElGamal & key generation, encryption, and decryption programs; exact
      encryption-path theorem; timed adapter & scheme PMFs and correctness \\
    \bottomrule
  \end{tabular}
\end{table}

\paragraph{One-time pad.}
The OTP parameter contains only the finite additive group and operations that
the construction uses.  There is no dummy scalar type or blanket scalar-action
instance.  \decl{keygenProgram}, \decl{encryptProgram}, and
\decl{decryptProgram} call typed sampling, addition, and subtraction
operations; their bounded variants attach certificates to those same programs.
Theorems such as \decl{scheme_encryptDist} expose the erased distribution, and
\decl{perfectOneTimeSecure} proves exact equality of the left and right
experiments before deriving \decl{oneTimeSecure} for every selected PPT
adversary.

\paragraph{DLog and DDH.}
\decl{CyclicAction} supplies a finite additive carrier, a finite scalar type,
an action, a generator, and generator surjectivity.  The stronger DDH parameter
\texttt{DecisionalCyclicAction} adds a commutative scalar monoid and the action
laws needed by DDH.  Its \decl{MulAction} is derived from the inherited action,
avoiding a second incoherent instance.  Setup produces the dependent public
parameter, after which result-indexed operations sample values of
\decl{pp.Scalar} and act on \decl{pp.Carrier}.  Program erasure lemmas connect
these exact paths to \decl{dLogProblem} and \decl{ddhProblem}.  Their
\decl{Assumption} definitions still quantify every PPT machine admitted by the
explicit adversary model and measure.

\paragraph{ElGamal.}
ElGamal reuses the DDH public parameter and handler rather than copying its
algebraic laws.  Encryption performs exactly one scalar sample, two scalar
actions, and one addition.  The theorem \decl{encryptProgram_exactCost}
decomposes every supported result into those four supported primitive results
and proves that the recorded cost is their ordered sum.  This theorem consults
the exact handler, not an upper-bound certificate.  Independently,
\decl{encryptProgram_costBound} certifies the budget and
\decl{encryptTimedMachine_runDist} proves that the machine adapter preserves
the encryption PMF.  The existing \decl{correct} theorem is stated over the
cost-erased scheme distributions.

The current library does not yet prove ElGamal IND-CPA security from the DDH
assumption.  The syntax and generic IND-CPA boundary exist, but presenting the
reduction as completed would overstate the mechanized result.
